\begin{lemma}
\label{lemma-colim-global}
Let $\mathcal{C}$ be a site. Let $S \subset \Ob(\Sh(\mathcal{C}))$
be a subset. Denote $*$ the final object of $\Sh(\mathcal{C})$. Assume
\begin{enumerate}
\item for some $K \in S$ the map $K \to *$ is surjective,
\item given a surjective map of sheaves $\mathcal{F} \to K$ with $K \in S$
there exists a $K' \in S$ and a map $K' \to \mathcal{F}$ such
that the composition $K' \to K$ is surjective,
\item given $K, K' \in S$ there is a surjection $K'' \to K \times K'$
with $K'' \in S$,
\item given $a, b : K \to K'$ with $K, K' \in S$ there exists a
surjection $K'' \to \text{Equalizer}(a, b)$ with $K'' \in S$, and
\item every $K \in S$ is quasi-compact
(Sites, Definition \ref{sites-definition-quasi-compact-topos}).
\end{enumerate}
Then for all $p \geq 0$ the map
$$
\colim_\lambda H^p(\mathcal{C}, \mathcal{F}_\lambda)
\longrightarrow
H^p(\mathcal{C}, \colim_\lambda \mathcal{F}_\lambda)
$$
is an isomorphism for every filtered diagram
$\Lambda \to \textit{Ab}(\mathcal{C})$, $\lambda \mapsto \mathcal{F}_\lambda$.
\end{lemma}

\begin{proof}
We will prove this by induction on $p$. The base case $p = 0$
follows from
Sites, Lemma \ref{sites-lemma-directed-colimits-global-sections} part (4).
We check the assumptions hold, but we urge the reader to skip this part.
Suppose $\mathcal{F} \to *$ is surjective.
Choose $K \in S$ and $K \to *$ surjective as in (1). Then
$\mathcal{F} \times K \to K$ is surjective. Choose
$K' \to \mathcal{F} \times K$ with $K' \in S$ and $K' \to K$
surjective as in (2).
Then there is a map $K' \to \mathcal{F}$ and $K' \to *$ is surjective.
Hence 
Sites, Lemma \ref{sites-lemma-directed-colimits-global-sections}
assumption (4)(a) is satisfied.
By Sites, Lemma \ref{sites-lemma-image-of-quasi-compact}, assumptions
(3) and (5) we see that $K \times K$ is quasi-compact
for all $K \in S$.
Hence 
Sites, Lemma \ref{sites-lemma-directed-colimits-global-sections}
assumption (4)(b) is satisfied.
This finishes the proof of the base case.

\medskip\noindent
Induction step. Assume the result holds for
$H^p$ for $p \leq p_0$ and for all topoi $\Sh(\mathcal{C})$
such that a set $S \subset \Ob(\Sh(\mathcal{C}))$ can be found
satisfying (1) -- (5). Arguing exactly as in the proof of
Lemma \ref{lemma-colim-works-over-collection}
we see that it suffices to show: given a filtered colimit
$\mathcal{I} = \colim \mathcal{I}_\lambda$ with
$\mathcal{I}_\lambda$ injective abelian sheaves, we have
$H^{p_0 + 1}(\mathcal{C}, \mathcal{I}) = 0$.
Choose $K \to *$ surjective with $K \in S$ as in (1).
Denote $K^n$ the $n$-fold self product of $K$.
Consider the spectral sequence
$$
E_1^{p, q} = H^q(K^{p + 1}, \mathcal{I}) \Rightarrow
H^{p + q}(*, \mathcal{I}) = H^{p + q}(\mathcal{C}, \mathcal{I})
$$
of Lemma \ref{lemma-cech-to-cohomology-sheaf-sets}. Recall that
$H^q(K^{p + 1}, \mathcal{F}) = H^q(\mathcal{C}/K^{p + 1}, j^{-1}\mathcal{F})$,
for any abelian sheaf on $\mathcal{C}$, see
Lemma \ref{lemma-cohomology-on-sheaf-sets}. We have
$j^{-1}\mathcal{I} = \colim j^{-1}\mathcal{I}_\lambda$
as $j^{-1}$ commutes with colimits. The restrictions
$j^{-1}\mathcal{I}_\lambda$ are injective abelian sheaves
on $\mathcal{C}/K^{p + 1}$ by Lemma \ref{lemma-cohomology-of-open}.
Below we will show that the induction hypothesis applies
to $\mathcal{C}/K^{p + 1}$ and hence we see that
$H^q(K^{p + 1}, \mathcal{I}) = \colim H^q(K^{p + 1}, \mathcal{I}_\lambda) = 0$
for $q < p_0 + 1$ (vanishing as $\mathcal{I}_\lambda$ is injective).
It follows that
$$
H^{p_0 + 1}(\mathcal{C}, \mathcal{I}) =
H^{p_0 + 1}\left(\ldots \to
H^0(K^{p_0}, \mathcal{I}) \to
H^0(K^{p_0 + 1}, \mathcal{I}) \to
H^0(K^{p_0 + 2}, \mathcal{I}) \to \ldots\right)
$$
Again using the induction hypothesis, the complex depicted on the right
hand side is the colimit over $\Lambda$ of the complexes
$$
\ldots \to
H^0(K^{p_0}, \mathcal{I}_\lambda) \to
H^0(K^{p_0 + 1}, \mathcal{I}_\lambda) \to
H^0(K^{p_0 + 2}, \mathcal{I}_\lambda) \to \ldots
$$
These complexes are exact as $\mathcal{I}_\lambda$ is an
injective abelian sheaf (follows from the spectral sequence
for example). Since filtered colimits are exact in the category
of abelian groups we obtain the desired vanishing.

\medskip\noindent
We still have to show that the induction hypothesis applies
to the site $\mathcal{C}/K^n$ for all $n \geq 1$. Recall
that $\Sh(\mathcal{C}/K^n) = \Sh(\mathcal{C})/K^n$, see
Sites, Lemma \ref{sites-lemma-localize-topos-site}.
Thus we may work in $\Sh(\mathcal{C})/K^n$.
Denote $S_n \subset \Ob(\Sh(\mathcal{C}/K^n)$ the set
of objects of the form $K' \to K^n$. We check each property in turn:
\begin{enumerate}
\item By (3) and induction there exists a surjection
$K' \to K^n$ with $K' \in S$. Then $(K' \to K^n) \to (K^n \to K^n)$
is a surjection in $\Sh(\mathcal{C})/K^n$ and $K^n \to K^n$
is the final object of $\Sh(\mathcal{C})/K^n$. Hence (1) holds for $S_n$,
\item Property (2) for $S_n$ is an immediate consequence
of (2) for $S$.
\item Let $a : K_1 \to K^n$ and $b : K_2 \to K^n$ be in $S_n$.
Then $(K_1 \to K^n) \times (K_2 \to K^n)$ is the object
$K_1 \times_{K^n} K_2 \to K^n$ of $\Sh(\mathcal{C})/K^n$.
The subsheaf $K_1 \times_{K^n} K_2 \subset K_1 \times K_2$ is the
equalizer of $a \circ \text{pr}_1$ and $b \circ \text{pr}_2$.
Write $a = (a_1, \ldots, a_n)$ and $b = (b_1, \ldots, b_n)$.
Pick $K_3 \to K_1 \times K_2$ surjective with $K_3 \in S$;
this is possibly by assumption (3) for $\mathcal{C}$.
Pick
$$
K_4
\longrightarrow
\text{Equalizer}(K_3 \to K_1 \times K_2 \xrightarrow{a_1, b_1} K)
$$
surjective with $K_4 \in S$.
This is possible by assumption (4) for $\mathcal{C}$.
Pick
$$
K_5
\longrightarrow
\text{Equalizer}(K_4 \to K_1 \times K_2 \xrightarrow{a_2, b_2} K)
$$
surjective with $K_5 \in S$.
Again this is possible. Continue in this fashion until we get
$$
K_{3 + n}
\longrightarrow
\text{Equalizer}(K_{2 + n} \to K_1 \times K_2 \xrightarrow{a_n, b_n} K)
$$
surjective with $K_{3 + n} \in S$. By construction
$K_{3 + n} \to K_1 \times_{K^n} K_2$
is surjective. Hence $(K_{3 + n} \to K^n)$ is in $S_n$ and
surjects onto the product $(K_1 \to K^n) \times (K_2 \to K^n)$.
Thus (3) holds for $S_n$.
\item Property (4) for $S_n$ is an immediate consequence of
property (4) for $S$.
\item Property (5) for $S_n$ is a consequence of
property (5) for $S$. Namely, an object $\mathcal{F} \to K^n$ of
$\Sh(\mathcal{C})/K^n$ corresponds to a quasi-compact object of
$\Sh(\mathcal{C}/K^n)$ if and only if $\mathcal{F}$ is a
quasi-compact object of $\Sh(\mathcal{C})$.
\end{enumerate}
This finishes the proof of the lemma.
\end{proof}